%
% Course IFJ @ FIT VUT Brno, 2015
% IFJ15 Interpreter Project
%
% Authors:
% Lukas Osadsky  - xosads00
% Pavol Plaskon  - xplask00
% Pavel Pospisil - xpospi88
% Matej Postolka - xposto02

\documentclass[a4paper, 12pt]{article}

% Použité nbalíčky
\usepackage[left=2cm,text={17cm, 24cm},top=3cm]{geometry}
\usepackage[czech]{babel}
\usepackage[utf8]{inputenc}
\usepackage[IL2]{fontenc}
%\usepackage{times}
\usepackage{graphics}
\usepackage{url}
% Požadavky automatu
\usepackage{pgf}
\usepackage{tikz}
\usetikzlibrary{arrows,automata}
% Požadavky tabulky
\usepackage{tabu}

% SMAZAT!!!!!!!
\usepackage{color}

\pagenumbering{arabic}
\providecommand{\uv}[1]{\quotedblbase #1\textquotedblleft}

\begin{document}
%%%%%%%%%%%%%%%%%%%%%%%%%%%%%%%%%%%%%%%%%%%%%%%%
%               Title page
%%%%%%%%%%%%%%%%%%%%%%%%%%%%%%%%%%%%%%%%%%%%%%%%
\begin{titlepage}

\begin{center}
\fontsize{25}{20}\selectfont{\textsc{Vysoké učení technické v~Brně}}\\
\vspace{\stretch{0.0075}}
\fontsize{21}{0}\textsc{\selectfont{Fakulta informačních technologií}}\\
\vspace{\stretch{0.15}}
%Logo
\begin{figure}[ht]
    \begin{center}
        \scalebox{1}{\includegraphics{logo_bw.eps}}
    \end{center}
\end{figure}

\vspace{\stretch{0.2}}
\LARGE{Dokumentace IFJ15}\\
\Large{Tým \textit{052}, varianta \textit{a/2/II}}
\vspace{\stretch{0.618}}
\end{center}

\begin{large}
\begin{tabbing}
    Další členové: \quad \= Postolka Matej \quad \= xposto02 \quad \= 25 \%\kill
    Vedoucí týmu:  \> Postolka Matěj \> xposto02 \> 25 \% \\
    Další členové: \> Osadský Lukáš  \> xosads00 \> 25 \% \\
             \> Plaskoň Pavol  \> xplask00 \> 25 \% \\
             \> Pospíšil Pavel \> xpospi88 \> 25 \% \\
\end{tabbing}
\end{large}

\end{titlepage}

%%%%%%%%%%%%%%%%%%%%%%%%%%%%%%%%%%%%%%%%%%%%%%%%
%                   Obsah
%%%%%%%%%%%%%%%%%%%%%%%%%%%%%%%%%%%%%%%%%%%%%%%%
\tableofcontents
\newpage
\textcolor{red}{Default font (tento) mi přijde krajšie, ještě se domluvíme}
%%%%%%%%%%%%%%%%%%%%%%%%%%%%%%%%%%%%%%%%%%%%%%%%
%                   Úvod
%%%%%%%%%%%%%%%%%%%%%%%%%%%%%%%%%%%%%%%%%%%%%%%%
\section{Úvod} \label{uvod}
Obecný pokec o stvoření světa
%%%%%%%%%%%%%%%%%%%%%%%%%%%%%%%%%%%%%%%%%%%%%%%%
%                   UNDER CONSTRUCTION
%%%%%%%%%%%%%%%%%%%%%%%%%%%%%%%%%%%%%%%%%%%%%%%%
\section{TODO} \label{todo}
\begin{verbatim}
kontrola čitelnosti kódu
kontrola komentovanosti kódů, zdroje! (např. odkud prvočísla do hash tabulky)
tvorba automatu http://madebyevan.com/fsm/
vývojový cyklus? speciální techniky
instrukční sadu?
Section "Implementace interpretu jazyka IFJ15" a pod lex, syn, sem, int?
Větší řádkování?

Zadání:
dokumentace.pdf      4-7 stran
1. strana            OK
Diagram automatu    "OK"
LL-gramatika        "OK"
prece tabulka       "OK"
POPIS ZPŮSOBU ŘEŠENÍ Z POHLEDU IFJ
    návrh
    implementace
    vývojový cyklus
    práce v týmu
    sleciální techniky
    algoritmy
Popis řadícího svého algoritmu
    řazení
    hledání v textu
    tabulky symbolů
Opět práce v týmu (kdo co jak)
Literatutra, reference, citace
!! Nevlastní materiál
!! Obecné popisy algoritmů

Pavel:
kontrola prece tabulky X změna za symboly < > = !
Kontrola stavů konečného automatu
Upravit automat pro čitelnost
\end{verbatim}
\newpage
%%%%%%%%%%%%%%%%%%%%%%%%%%%%%%%%%%%%%%%%%%%%%%%%
%                   Práce v týmu
%%%%%%%%%%%%%%%%%%%%%%%%%%%%%%%%%%%%%%%%%%%%%%%%
\section{Práce v týmu} \label{team}

%%%%%%%%%%%%%%%%%%%%%%%%%%%%%%%%%%%%%%%%%%%%%%%%
%                   Lexikální analyzátor
%%%%%%%%%%%%%%%%%%%%%%%%%%%%%%%%%%%%%%%%%%%%%%%%
\section{Lexikální analyzátor} \label{lexer}
\textcolor{red}{Lexikální analýza?}

Lexikální analyzátor je vstupní a nejjednodušší \textcolor{red}{Nezdůroazňoval bych} část překladače. Je založen na deterministickém \textcolor{red}{bejvávalo} konečném automatu,
jehož hlavním úkolem je čtení zdrojového souboru a na základě lexikálních pravidel jazyka rozdělit jednotlivé znaky \textcolor{red}{množiny znaků}
souboru na lexikální části (lexémy). Rozpoznané lexémy jsou reprezenované strukturou token, která obsahuje informace o typu
tokenu a jeho data. Jeho vedlejší úlohou je odstraňování všech komentářů a bílých znaků, které \textcolor{red}{neboť} nejsou potřebné pro
následné zpracování. Princip fungování lexikálního analyzátoru reprezuntuje příloha č.1 \textcolor{red}{A}, ve které je zobrazeno jeho schéma.
Činnost lexikálního analyzátoru je přímo řízena syntaktickým analyzátorem. Jeho výstupem \textcolor{red}{čeho?} je token,
který je zároveň vstupem do syntaktické analýzy.

%%%%%%%%%%%%%%%%%%%%%%%%%%%%%%%%%%%%%%%%%%%%%%%%
%                   Syntaktický analyzátor
%%%%%%%%%%%%%%%%%%%%%%%%%%%%%%%%%%%%%%%%%%%%%%%%
\section{Syntaktický analyzátor} \label{parser}

\textcolor{red}{Syntaktická a sémantická analýza?}

\subsection{Zpracování výrazů}
\textcolor{red}{A FUNKCÍ! Má jediná chvíle slávy?}

%%%%%%%%%%%%%%%%%%%%%%%%%%%%%%%%%%%%%%%%%%%%%%%%
%                   Vestavěné funkce
%%%%%%%%%%%%%%%%%%%%%%%%%%%%%%%%%%%%%%%%%%%%%%%%
\section{Vestavěné funkce????} \label{vestavene}

%%%%%%%%%%%%%%%%%%%%%%%%%%%%%%%%%%%%%%%%%%%%%%%%
%                   Interpret
%%%%%%%%%%%%%%%%%%%%%%%%%%%%%%%%%%%%%%%%%%%%%%%%
\section{Interpret} \label{interpret}

\subsection{Řadíci algoritmus -- Heap Sort}
Funkce pro seřazení prvků v poli.

\subsection{Vyhledávaní podřetězce -- Knuth-Morris-Pratt}
Vyhledání podřetězce v řetězci ve vestavěné funkci \texttt{find} je řešeno
algoritmem Knuth-Morris-Pratt. Základem algoritmu je vytvoření masky, tzv.
\texttt{Fail vector}. Je to \textcolor{red}{nevábný počátek věty} pole celých čísel o délce hledaného textu. Ke
každému písmenu hledaného řetězce je přiřazeno číslo, které určuje index, kam \textcolor{red}{pro návrat programu}
se má program vrátit v případě neshody znaků.

\subsection{Tabulka s rozptýlenými položkami}
Datová struktura použitá pro tabulky \textcolor{red}{tabulku?} symbolů. \textcolor{red}{Její} Výhodou je rychlost vyhledávání
položek. Základem je pole ukazatelů na
jednotlivé položky. Položky obsahují svůj klíč, data a ukazatel na další
položku, aby mohly být propojené v jednosměrně vázaný lineární
seznam \textcolor{red}{Ajajaj, to by se mi asi nemělo stávat. Utíkající odstavec} (seznam synonym) \textcolor{red}{-- seznam synonym}. V případě ideální hashovací funkce není propojení v
seznam potřebné, \textcolor{red}{a} čas přístupu k položkám konstantní. Nalezení takové
funkce je ale problematické \textcolor{red}{není triviální}. V případě konfliktů \textcolor{red}{konfliktu?} se čas nalezení položky
prodloužuje o \textcolor{red}{dobu?} prohledání lineárního seznamu.

\newpage
%%%%%%%%%%%%%%%%%%%%%%%%%%%%%%%%%%%%%%%%%%%%%%%%
%                   Přílohy
%%%%%%%%%%%%%%%%%%%%%%%%%%%%%%%%%%%%%%%%%%%%%%%%
\appendix
% AUTOMAT
\section{Diagram konečného automatu lexikální analýzy} \label{automat}
\begin{center}
\begin{tikzpicture}[scale=0.165]
\small
\tikzstyle{every node}+=[inner sep=0pt]
\draw [black] (40,-55) circle (3);
\draw (40,-55) node {\em start};
\draw [black] (7,-114.5) circle (3);
\draw (7,-114.5) node {\em +-*,;};
\draw [black] (7,-114.5) circle (2.4);
\draw [black] (11.7,-131.3) circle (3);
\draw (11.7,-131.3) node {\em (){}};
\draw [black] (11.7,-131.3) circle (2.4);
\draw [black] (7,-87.3) circle (3);
\draw (7,-87.3) node {\em =};
\draw [black] (7,-87.3) circle (2.4);
\draw [black] (77,-51.2) circle (3);
\draw (77,-51.2) node {\em /};
\draw [black] (77,-51.2) circle (2.4);
\draw [black] (94.1,-43.4) circle (3);
\draw (94.1,-43.4) node {\em lbcom};
\draw [black] (91.7,-63.2) circle (3);
\draw (91.7,-63.2) node {\em lcom};
\draw [black] (69.1,-39.7) circle (3);
\draw (69.1,-39.7) node {\em rbcom};
\draw [black] (76,-24) circle (3);
\draw (76,-24) node {\em <};
\draw [black] (76,-24) circle (2.4);
\draw [black] (47.8,-19.9) circle (3);
\draw (47.8,-19.9) node {\em >};
\draw [black] (47.8,-19.9) circle (2.4);
\draw [black] (77,-5.6) circle (3);
\draw (77,-5.6) node {\em cin};
\draw [black] (77,-5.6) circle (2.4);
\draw [black] (94.1,-21.3) circle (3);
\draw (94.1,-21.3) node {\em leseq};
\draw [black] (94.1,-21.3) circle (2.4);
\draw [black] (37.7,-9.6) circle (3);
\draw (37.7,-9.6) node {\em cout};
\draw [black] (37.7,-9.6) circle (2.4);
\draw [black] (62.2,-11.7) circle (3);
\draw (62.2,-11.7) node {\em greq};
\draw [black] (62.2,-11.7) circle (2.4);
\draw [black] (6.2,-105) circle (3);
\draw (6.2,-105) node {\em eq};
\draw [black] (6.2,-105) circle (2.4);
\draw [black] (10.5,-64.6) circle (3);
\draw (10.5,-64.6) node {\em !};
\draw [black] (6.2,-54.3) circle (3);
\draw (6.2,-54.3) node {\em neq};
\draw [black] (6.2,-54.3) circle (2.4);
\draw [black] (10.5,-76.3) circle (3);
\draw (10.5,-76.3) node {\em eof};
\draw [black] (10.5,-76.3) circle (2.4);
\draw [black] (95.9,-86.5) circle (3);
\draw (95.9,-86.5) node {\em quot};
\draw [black] (73.2,-82.8) circle (3);
\draw (73.2,-82.8) node {\em str};
\draw [black] (73.2,-82.8) circle (2.4);
\draw [black] (76,-107.1) circle (3);
\draw (76,-107.1) node {\em esc};
\draw [black] (101.2,-125.5) circle (3);
\draw (101.2,-125.5) node {\em hex};
\draw [black] (58.5,-82.8) circle (3);
\draw (58.5,-82.8) node {\em int};
\draw [black] (58.5,-82.8) circle (2.4);
\draw [black] (44.9,-98.5) circle (3);
\draw (44.9,-98.5) node {\em expe};
\draw [black] (69.1,-99.5) circle (3);
\draw (69.1,-99.5) node {\em int.};
\draw [black] (17,-19.9) circle (3);
\draw (17,-19.9) node {\em id};
\draw [black] (17,-19.9) circle (2.4);
\draw [black] (23.1,-7.9) circle (3);
\draw (23.1,-7.9) node {\em keyword};
\draw [black] (23.1,-7.9) circle (2.4);
\draw [black] (80.4,-133.6) circle (3);
\draw (80.4,-133.6) node {\em double};
\draw [black] (80.4,-133.6) circle (2.4);
\draw [black] (49,-137.1) circle (3);
\draw (49,-137.1) node {\em exp};
\draw [black] (49,-137.1) circle (2.4);
\draw [black] (29.3,-135.3) circle (3);
\draw (29.3,-135.3) node {\em exsqn};
\draw [black] (37.02,-54.772) arc (293.36193:5.36193:2.25);
\draw (33.55,-50.31) node [left] {\em white\mbox{ }space};
\fill [black] (38.37,-52.5) -- (38.51,-51.56) -- (37.59,-51.96);
\draw [black] (19.7,-55) -- (37,-55);
\fill [black] (37,-55) -- (36.2,-54.5) -- (36.2,-55.5);
\draw [black] (39.093,-57.859) arc (-19.46398:-71.76436:46.378);
\fill [black] (9.88,-86.45) -- (10.79,-86.68) -- (10.48,-85.73);
\draw (28.89,-76.03) node [below] {\em =};
\draw [black] (39.263,-57.908) arc (-14.89614:-43.13128:127.582);
\fill [black] (9.08,-112.33) -- (9.99,-112.09) -- (9.26,-111.41);
\draw (28.21,-88.19) node [right] {\em +*-,;};
\draw [black] (40.161,-57.996) arc (2.21629:-42.91636:98.858);
\fill [black] (13.78,-129.13) -- (14.69,-128.89) -- (13.95,-128.21);
\draw (34.82,-97) node [right] {\em (){}};
\draw [black] (42.702,-53.697) arc (114.04706:77.68068:50.557);
\fill [black] (74.09,-50.47) -- (73.41,-49.81) -- (73.2,-50.79);
\draw (57.98,-49.01) node [above] {\em /};
\draw [black] (79.408,-49.413) arc (124.31794:104.72135:37.986);
\fill [black] (91.17,-44.05) -- (90.27,-43.77) -- (90.52,-44.73);
\draw (84.08,-45.72) node [above] {\em *};
\draw [black] (79.72,-52.464) arc (62.43568:39.11297:32.564);
\fill [black] (89.92,-60.79) -- (89.8,-59.85) -- (89.02,-60.48);
\draw (86.03,-55.61) node [above] {\em /};
\draw [black] (92.235,-66.14) arc (38.0546:-249.9454:2.25);
\draw (88.19,-70.3) node [below] {\em else};
\fill [black] (89.69,-65.41) -- (88.75,-65.51) -- (89.37,-66.3);
\draw [black] (95.133,-40.596) arc (187.51736:-100.48264:2.25);
\draw (99.52,-37.28) node [right] {\em else};
\fill [black] (96.95,-42.51) -- (97.81,-42.9) -- (97.68,-41.91);
\draw [black] (70.517,-37.063) arc (145.12644:18.03624:12.98);
\fill [black] (70.52,-37.06) -- (71.38,-36.69) -- (70.56,-36.12);
\draw (83.43,-31.05) node [above] {\em *};
\draw [black] (72.099,-39.642) arc (89.82687:73.3358:67.48);
\fill [black] (91.25,-42.48) -- (90.62,-41.77) -- (90.34,-42.73);
\draw (82.53,-39.68) node [above] {\em else};
\draw [black] (42.993,-55.201) arc (85.83691:76.13817:274.264);
\fill [black] (42.99,-55.2) -- (43.75,-55.76) -- (43.83,-54.76);
\draw (66.57,-57.24) node [above] {\em \textbackslash n};
\draw [black] (42.365,-53.155) arc (126.94089:108.52749:84.291);
\fill [black] (42.37,-53.15) -- (43.31,-53.07) -- (42.7,-52.27);
\draw (53.02,-45.42) node [above] {\em /};
\draw [black] (41.643,-52.49) arc (145.72712:115.7371:80.666);
\fill [black] (73.27,-25.25) -- (72.34,-25.15) -- (72.77,-26.05);
\draw (54.59,-36.3) node [above] {\em <};
\draw [black] (47.785,-22.9) arc (-1.39065:-23.66697:77.89);
\fill [black] (47.78,-22.9) -- (47.27,-23.69) -- (48.27,-23.71);
\draw (46.7,-38.28) node [right] {\em >};
\draw [black] (75.333,-21.077) arc (-170.18655:-196.03513:28.304);
\fill [black] (76.02,-8.43) -- (75.32,-9.06) -- (76.28,-9.34);
\draw (74.38,-14.69) node [left] {\em <};
\draw [black] (78.777,-22.869) arc (109.52567:87.44298:32.562);
\fill [black] (91.11,-21.03) -- (90.34,-20.49) -- (90.29,-21.49);
\draw (84.46,-20.75) node [above] {\em =};
\draw [black] (45.096,-18.61) arc (-120.5542:-150.56921:16.987);
\fill [black] (38.94,-12.33) -- (38.89,-13.27) -- (39.76,-12.78);
\draw (41.08,-17.35) node [left] {\em >};
\draw [black] (49.971,-17.831) arc (130.58707:108.73115:28.355);
\fill [black] (59.31,-12.51) -- (58.4,-12.29) -- (58.72,-13.24);
\draw (53.33,-14.22) node [above] {\em =};
\draw [black] (7.643,-90.229) arc (9.39106:-14.56682:28.726);
\fill [black] (7.1,-102.14) -- (7.79,-101.49) -- (6.82,-101.24);
\draw (8.57,-96.23) node [right] {\em =};
\draw [black] (37.384,-56.467) arc (-61.91096:-82.03681:71.936);
\fill [black] (13.48,-64.25) -- (14.34,-64.63) -- (14.2,-63.64);
\draw (26.51,-61.94) node [below] {\em !};
\draw [black] (8.605,-62.281) arc (-146.63075:-168.05055:14.588);
\fill [black] (6.52,-57.28) -- (6.19,-58.16) -- (7.17,-57.96);
\draw (6.59,-60.81) node [left] {\em =};
\draw [black] (38.338,-57.497) arc (-35.44861:-72.89027:47.931);
\fill [black] (13.39,-75.51) -- (14.3,-75.75) -- (14.01,-74.79);
\draw (29.13,-69.06) node [below] {\em eof};
\draw [black] (42.832,-55.989) arc (70.25846:50.93848:173.594);
\fill [black] (93.59,-84.59) -- (93.28,-83.7) -- (92.65,-84.47);
\draw (70.32,-67.64) node [above] {\em ''};
\draw [black] (95.246,-83.584) arc (220.37178:-67.62822:2.25);
\draw (98.93,-79.29) node [above] {\em else};
\fill [black] (97.82,-84.21) -- (98.75,-84.07) -- (98.1,-83.31);
\draw [black] (92.983,-87.194) arc (-79.88974:-118.62536:26.331);
\fill [black] (75.75,-84.38) -- (76.21,-85.21) -- (76.69,-84.33);
\draw (83.74,-87.84) node [below] {\em ''};
\draw [black] (95.237,-89.424) arc (-16.18725:-71.83235:25.191);
\fill [black] (78.9,-106.34) -- (79.82,-106.56) -- (79.5,-105.61);
\draw (89.69,-101.37) node [right] {\em /};
\draw [black] (77.714,-104.638) arc (144.09148:127.88892:80.608);
\fill [black] (93.5,-88.3) -- (92.56,-88.39) -- (93.17,-89.18);
\draw (84.5,-94.44) node [left] {\em n t \textbackslash ''};
\draw [black] (101.606,-128.461) arc (35.54413:-252.45587:2.25);
\draw (96.39,-132.51) node [below] {\em isdigit};
\fill [black] (99.1,-127.62) -- (98.16,-127.68) -- (98.74,-128.5);
\draw [black] (100.367,-122.618) arc (-164.56024:-179.96188:124.712);
\fill [black] (95.87,-89.5) -- (95.37,-90.3) -- (96.37,-90.3);
\draw (96.32,-106.37) node [left] {\em i=2};
\draw [black] (78.779,-108.23) arc (66.39614:41.33311:58.489);
\fill [black] (99.28,-123.2) -- (99.12,-122.27) -- (98.37,-122.93);
\draw (90.85,-114.09) node [above] {\em x};
\draw [black] (42.397,-56.803) arc (51.17912:16.10584:46.023);
\fill [black] (57.76,-79.89) -- (58.02,-78.99) -- (57.06,-79.26);
\draw (52.47,-65.83) node [right] {\em isdigit};
\draw [black] (55.522,-83.049) arc (302.51909:14.51909:2.25);
\draw (51.47,-79.46) node [left] {\em isdigit};
\fill [black] (56.49,-80.59) -- (56.48,-79.64) -- (55.64,-80.18);
\draw [black] (56.54,-85.07) -- (46.86,-96.23);
\fill [black] (46.86,-96.23) -- (47.77,-95.96) -- (47.01,-95.3);
\draw (51.15,-89.2) node [left] {\em e\mbox{ }E};
\draw [black] (60.87,-84.636) arc (48.7983:16.01079:25.048);
\fill [black] (68.45,-96.57) -- (68.71,-95.67) -- (67.75,-95.94);
\draw (66.14,-88.76) node [right] {\em .};
\draw [black] (19.42,-21.673) arc (52.23036:14.24094:55.835);
\fill [black] (19.42,-21.67) -- (19.75,-22.56) -- (20.36,-21.77);
\draw (32.54,-33.88) node [right] {\em \_ isalpha};
\draw [black] (15.339,-22.384) arc (-6.02889:-294.02889:2.25);
\draw (10.5,-24.5) node [left] {\em \_ isalnum};
\fill [black] (14.02,-20.09) -- (13.27,-19.51) -- (13.17,-20.51);
\draw [black] (17.244,-16.916) arc (-190.7275:-223.164:14.178);
\fill [black] (20.83,-9.86) -- (19.92,-10.1) -- (20.65,-10.78);
\draw (17.84,-12.01) node [left] {\em iskeyword};
\draw [black] (78.395,-131.369) arc (-140.13542:-183.19643:41.454);
\fill [black] (78.4,-131.37) -- (78.27,-130.43) -- (77.5,-131.08);
\draw (70.1,-118.53) node [left] {\em isdigit};
\draw [black] (81.377,-136.424) arc (46.82314:-241.17686:2.25);
\draw (78.71,-141.07) node [below] {\em isdigit};
\fill [black] (78.75,-136.09) -- (77.84,-136.33) -- (78.57,-137.02);
\draw [black] (77.643,-132.417) arc (-114.55876:-154.79201:64.458);
\fill [black] (46.11,-101.24) -- (46,-102.18) -- (46.91,-101.75);
\draw (57.29,-120.1) node [below] {\em e\mbox{ }E};
\draw [black] (48.28,-134.188) arc (-166.76078:-181.11304:131.577);
\fill [black] (48.28,-134.19) -- (48.58,-133.29) -- (47.61,-133.52);
\draw (44.87,-118.05) node [left] {\em isdigit};
\draw [black] (50.915,-134.806) arc (167.87147:-120.12853:2.25);
\draw (55.9,-133.38) node [right] {\em isdigit};
\fill [black] (51.99,-137.22) -- (52.66,-137.88) -- (52.87,-136.9);
\draw [black] (29.923,-132.366) arc (167.12944:146.92511:97.148);
\fill [black] (29.92,-132.37) -- (30.59,-131.7) -- (29.61,-131.47);
\draw (34.45,-115.15) node [left] {\em +-};
\draw [black] (32.273,-134.907) arc (95.27929:74.27941:38.224);
\fill [black] (46.15,-136.17) -- (45.51,-135.48) -- (45.24,-136.44);
\draw (39.75,-134.17) node [above] {\em isdigit};
\end{tikzpicture}
\end{center}
\newpage

% LL-GRAMATIKA
\section{LL--gramatika} \label{llgram}
\textcolor{red}{Matěj je rebel! Víc řádků než by snesla jedna strana.} \\
\textcolor{red}{Nesázet verbatimem?}
{\scriptsize
\begin{verbatim}
PROG->FUNCTION_DECL PROG
PROG->eps
FUNCTION_DECL->DATA_TYPE t_identifier t_lround_bracket FUNC_DECL_PARAMS t_rround_bracket NESTED_BLOCK
DATA_TYPE->t_int
DATA_TYPE->t_double
DATA_TYPE->t_string
FUNC_DECL_PARAMS->DATA_TYPE t_identifier FUNC_DECL_PARAMS_NEXT
FUNC_DECL_PARAMS_NEXT->t_comma FUNC_DECL_PARAMS
FUNC_DECL_PARAMS->eps
FUNC_DECL_PARAMS_NEXT->eps
NESTED_BLOCK->t_lcurly_bracket NBC t_rcurly_bracket
NBC->DECL_OR_ASSIGN NBC
DECL_OR_ASSIGN->DATA_TYPE t_identifier DECL_ASSIGN t_semicolon
DECL_OR_ASSIGN->t_auto t_identifier t_assign EXPRESSION t_semicolon
DECL_ASSIGN->t_assign EXPRESSION
DECL_ASSIGN->eps
NBC->FCALL_OR_ASSIGN NBC
FCALL_OR_ASSIGN->t_identifier FOA_PART2
FOA_PART2->t_lround_bracket FUNCTION_CALL_PARAMS t_rround_bracket t_semicolon
FOA_PART2->t_assign EXPRESSION t_semicolon
HARD_VALUE->t_int_value
HARD_VALUE->t_double_value
HARD_VALUE->t_string_value
FUNCTION_CALL_PARAMS->FUNCTION_CALL_PARAM FUNCTION_CALL_PARAMS_NEXT
FUNCTION_CALL_PARAMS->eps
FUNCTION_CALL_PARAM->t_identifier
FUNCTION_CALL_PARAM->HARD_VALUE
FUNCTION_CALL_PARAMS_NEXT->t_comma FUNCTION_CALL_PARAMS
FUNCTION_CALL_PARAMS_NEXT->eps
NBC->BUILTIN_CALL NBC
BUILTIN_CALL->BUILTIN_FUNC t_lround_bracket FUNCTION_CALL_PARAMS t_rround_bracket t_semicolon
BUILTIN_FUNC->token_bf_length
BUILTIN_FUNC->token_bf_substr
BUILTIN_FUNC->token_bf_concat
BUILTIN_FUNC->token_bf_find
BUILTIN_FUNC->token_bf_sort
NBC->IF_STATEMENT NBC
IF_STATEMENT->t_if t_lround_bracket EXPRESSION t_rround_bracket NESTED_BLOCK ELSE_STATEMENT
ELSE_STATEMENT->t_else NESTED_BLOCK
ELSE_STATEMENT->eps
NBC->COUT NBC
COUT->t_cout t_cout_bracket COUT_OUTPUT COUT_NEXT t_semicolon
COUT_OUTPUT->t_identifier
COUT_OUTPUT->HARD_VALUE
COUT_NEXT->t_cout_bracket COUT_OUTPUT COUT_NEXT
COUT_NEXT->eps
NBC->CIN NBC
CIN->t_cin t_cin_bracket t_identifier CIN_NEXT t_semicolon
CIN_NEXT->t_cin_bracket t_identifier CIN_NEXT
CIN_NEXT->eps
NBC->FOR_STATEMENT NBC
FOR_STATEMENT->t_for t_lround_bracket FOR_DECLARATION FOR_EXPR FOR_ASSIGN t_rround_bracket NESTED_BLOCK
FOR_DECLARATION->DATA_TYPE t_identifier DECL_ASSIGN t_semicolon
FOR_DECLARATION->t_auto t_identifier t_assign EXPRESSION t_semicolon
FOR_EXPR->EXPRESSION t_semicolon
FOR_ASSIGN->t_identifier t_assign EXPRESSION
NBC->NESTED_BLOCK NBC
NBC->RETURN
RETURN->t_return EXPRESSION t_semicolon
NBC->eps
\end{verbatim}
}
\newpage

% PRECEDENČNÍ TABULKA
\section{Precedenční tabulka} \label{prece}
\textcolor{red}{záhlaví tabulky??}
\begin{center}
\setlength{\extrarowsep}{3pt}
\begin{tabu}{|l|[2pt]l|l|l|l|l|l|l|l|l|l|l|l|l|}
\hline
st\textbackslash in              & +  & -  & *  & /  & (  & )  & id & \textless & \textgreater & \textless= & \textgreater= & == & != \\ \tabucline[2pt]{-}
+             & HI & HI & LO & LO & LO & HI & LO & HI        & HI           & HI         & HI            & HI & HI \\ \hline
-             & HI & HI & LO & LO & LO & HI & LO & HI        & HI           & HI         & HI            & HI & HI \\ \hline
*             & HI & HI & HI & HI & LO & HI & LO & HI        & HI           & HI         & HI            & HI & HI \\ \hline
/             & HI & HI & HI & HI & LO & HI & LO & HI        & HI           & HI         & HI            & HI & HI \\ \hline
(             & LO & LO & LO & LO & LO & EQ & LO & LO        & LO           & LO         & LO            & LO & LO \\ \hline
)             & HI & HI & HI & HI & ER & HI & ER & HI        & HI           & HI         & HI            & HI & HI \\ \hline
id            & HI & HI & HI & HI & ER & HI & ER & HI        & HI           & HI         & HI            & HI & HI \\ \hline
\textless     & LO & LO & LO & LO & LO & HI & LO & HI        & HI           & HI         & HI            & HI & HI \\ \hline
\textgreater  & LO & LO & LO & LO & LO & HI & LO & HI        & HI           & HI         & HI            & HI & HI \\ \hline
\textless=    & LO & LO & LO & LO & LO & HI & LO & HI        & HI           & HI         & HI            & HI & HI \\ \hline
\textgreater= & LO & LO & LO & LO & LO & HI & LO & HI        & HI           & HI         & HI            & HI & HI \\ \hline
==            & LO & LO & LO & LO & LO & HI & LO & LO        & LO           & LO         & LO            & HI & HI \\ \hline
!=            & LO & LO & LO & LO & LO & HI & LO & LO        & LO           & LO         & LO            & HI & HI \\ \hline
\end{tabu}
\end{center}
\newpage

% INSTRUKČNÍ SADA TROJADRESNÉHO KÓDU
\section{Instrukční sada trojadresného kódu} \label{instrukce}
\textcolor{red}{Přidat popisky a nesázet přes Verbatim}
\begin{verbatim}
    INS_ASSIGN
    INS_ADD
    INS_SUB
    INS_MUL
    INS_DIV
    INS_EQ
    INS_NEQ
    INS_GREATER
    INS_GREATEQ
    INS_LESSER
    INS_LESSEQ
    INS_JMP
    INS_CJMP
    INS_LAB
    INS_PUSH
    INS_CALL
    INS_RET
    INS_PUSH_TAB
    INS_POP_TAB
    INS_LENGTH
    INS_SUBSTR
    INS_CONCAT
    INS_FIND
    INS_SORT
    INS_CIN
    INS_COUT
\end{verbatim}

\end{document}

